\section{Discussion and Future Work}
\label{sec:discussion}

\myparagraph{Generalizability to LoRA Variants}
Our kernel fusion design is extensible to other popular LoRA variants like DoRA~\cite{liu2024dora} and VeRA~\cite{kopiczko2023vera}.
These methods typically add pre- or post-processing functions around the core LoRA computation.
Our optimizations are orthogonal to these modifications, and users can define prologue/epilogue functions to extend our kernels.
While manual extension is effective, a more general approach is to integrate our fusion patterns into a compiler framework.
As future work, we plan to leverage \texttt{torch.compile} by adding compiler annotations as hints that guide the fusion process of the LoRA pattern.
This would automate the optimization for both existing and future LoRA variants, eliminating the need for manual kernel development and system expertise from users.

\myparagraph{Generalizability to Quantization}
The kernels proposed in \Name can be directly applied to 4-bit QLoRA~\cite{dettmers2023qlora}. 
Current QLoRA implementations dequantize 4-bit weights to half-precision before LoRA computation, allowing our kernels to work without modification.
While dequantization could be fused with the LoRA path, recent work shows that two-step approaches are often more performant for large token counts~\cite{ding2025tilus}.
